% !TeX program = xelatex
% ============================================================================
%  第 14 课　复习、总结与展望　—— 学生版讲义
%  与 02-课件/第14课-复习总结与展望/第14课-复习总结与展望.tex 同步
% ============================================================================

\xhlesson{14}{复习、总结与展望}{}

\begin{mubiao}
  \begin{itemize}
    \item 能自己说出这门课学过的三类「绕圈」。
    \item 能把它们和群的四条规则对上。
    \item 知道群之后还有环和域，知道五次方程的故事。
  \end{itemize}
\end{mubiao}

\section*{一、把 14 节课串成一条线}

\noindent 今天不学新东西。老师会在黑板中间写一个大字「群」，往三个方向拉线，
你要做的是一边听，一边把这张图补到下面。

\begin{dongshou}[画出主线图]
  在下面空白处画出三条线，每条线上写两个你记得住的例子。

  \liukong{7}
\end{dongshou}

\begin{sikao}
  三条线为什么会汇到一起？它们共同满足的，是哪四条规则？\\
  提示：想想「运算的结果还在不在里面」「有没有什么都不做的元」「能不能撤回一步」「两次三次怎么合并」。
  \liukong{4}
\end{sikao}

\begin{definition}[群]
  一个集合配上一种运算，如果满足下面四条，就叫\textbf{群}：
  \begin{enumerate}
    \item \textbf{封闭}：任取 \(a,b\)，\(a*b\) 还在这里面；
    \item \textbf{结合}：\((a*b)*c=a*(b*c)\)；
    \item \textbf{单位元}：存在 \(e\)，使 \(a*e=e*a=a\)；
    \item \textbf{逆元}：每个 \(a\) 都有 \(b\)，使 \(a*b=b*a=e\)。
  \end{enumerate}
\end{definition}

\section*{二、抽查}

\begin{dongshou}[合上书，把答案写在横线上]
  \begin{enumerate}
    \item 模 \(12\) 加法里，\(5\) 的逆元是 \underline{\hspace{3em}}。
    \item 正三角形有 \underline{\hspace{2em}} 个对称动作；正方形有 \underline{\hspace{2em}} 个。
    \item 三角形的对称里，\(rf\) 和 \(fr\) 一样吗？\underline{\hspace{4em}}
    \item \(2^{100}\) 除以 \(7\) 余 \underline{\hspace{2em}}。
  \end{enumerate}
\end{dongshou}

\begin{sikao}
  上面四个答案，各属于三条主线里的哪一支？在每条答案后面写一个「左 / 右 / 下」。
  \liukong{2}
\end{sikao}

\section*{三、前面还有什么}

\noindent 分配律你们早就在用了：\(3\times(4+5)=3\times4+3\times5\)。

\begin{definition}[环]
  一个集合上同时有加法和乘法，如果下面三条都满足，就叫做\textbf{环}：
  \begin{enumerate}
    \item \textbf{加法}：满足第 4 课那四条规则（封闭、结合、单位元、逆元），而且加法可以交换：\(a+b=b+a\)；
    \item \textbf{乘法}：封闭、结合、\textbf{有单位元 \(1\)}（\(a\cdot1=1\cdot a=a\)）、而且乘法可以交换：\(a\cdot b=b\cdot a\)——但乘法不要求每个数都有逆元；
    \item \textbf{乘法对加法有分配律}：\(a\cdot(b+c)=a\cdot b+a\cdot c\)，\((b+c)\cdot a=b\cdot a+c\cdot a\)。
  \end{enumerate}
  乘法可以交换，所以它又叫\textbf{交换环}——乘法不交换的环也有，这门课不去碰它。
\end{definition}

\begin{definition}[域]
  在上面这个环的基础上，如果每个非零元素都有乘法逆元（能做除法），就叫做\textbf{域}。
\end{definition}

\begin{example}
  整数 \(\mathbb{Z}\) 是环但不是域（\(2\) 没有整数逆元，\(\tfrac12\notin\mathbb{Z}\)）；有理数 \(\mathbb{Q}\) 是域。
\end{example}

\begin{sikao}
  为什么说「整数是环，不是域」？在整数里举一个除不了法的例子。
  \liukong{2}
\end{sikao}

\section*{四、附：一、二、三、四次方程的公式长什么样（只看个样子，不用记）}

\noindent 下面这几页\textbf{只看个样子，不用记、不用算}。
（\(\sqrt{\ }\) 是平方根，\(\sqrt[3]{\ }\) 是立方根；这里只要认得这个记号就行。）

\begin{example}[一元二次方程：配方就能解]
  \[ ax^2+bx+c=0\quad(a\neq0) \]
  两边除以 \(a\)，再令中间变量
  \[ t=x+\frac{b}{2a}, \]
  就化成
  \[ t^2=\frac{b^2-4ac}{4a^2}. \]
  令 \(\Delta=b^2-4ac\)（叫\textbf{判别式}），于是
  \[ t=\pm\frac{\sqrt{\Delta}}{2a},\qquad x=\frac{-b\pm\sqrt{b^2-4ac}}{2a}. \]
  \textbf{一句话}：先把 \(x\) 平移到 \(t\)，方程就只剩一个平方。
  （\(\Delta>0\) 两个根，\(\Delta=0\) 一个根，\(\Delta<0\) 没有实数根。）

  真正把这几个公式一条条全写出来，能占满一整块黑板；所以数学家后来才问：
  五次能不能也有这样一条路？
  \liukong{1}
\end{example}

\begin{example}[一元三次方程：卡尔达诺设 \(t=u+v\)]
  \[ t^3+pt+q=0 \]
  （任何三次方程都能先平移消掉平方项，变成这个「缺项」的样子。）
  令
  \[ t=u+v,\qquad\text{并约定}\ uv=-\frac{p}{3}, \]
  代入后方程变成
  \[ u^3+v^3=-q, \]
  而由约定又有
  \[ u^3v^3=-\frac{p^3}{27}. \]
  所以 \(u^3\) 与 \(v^3\) 是一元二次方程
  \[ z^2+qz-\frac{p^3}{27}=0 \]
  的两个根。\textbf{解出这个二次方程，就得到 \(u^3,v^3\)，再开立方就得到 \(t=u+v\)。}

  \textbf{一句话}：三次方程被「两个中间变量」降成了一个二次方程。

  真正把这几个公式一条条全写出来，能占满一整块黑板；所以数学家后来才问：
  五次能不能也有这样一条路？
  \liukong{1}
\end{example}

\begin{example}[一元四次方程：费拉里配两次方]
  \[ t^4+pt^2+qt+r=0 \]
  （同样先平移消掉三次项。）第一步配方：
  \[ \left(t^2+\frac{p}{2}\right)^2=-\frac{p^2}{4}+r-qt. \]
  第二步，\textbf{引入中间变量 \(y\)}，给左边添上 \(2y\left(t^2+\frac p2\right)+y^2\)，
  让它仍然是完全平方：
  \[ \left(t^2+\frac{p}{2}+y\right)^2
     =2y\,t^2-qt+\left(y^2+py-\frac{p^2}{4}+r\right). \]
  只要右边的判别式为 \(0\)（即 \(q^2=8y^3+8py^2-8ry\)，这是\textbf{一个关于 \(y\) 的三次方程}），
  右边也是完全平方；于是两边开方，就得到\textbf{两个二次方程}。

  \textbf{一句话}：四次方程靠「中间变量 \(y\)」拆成两个二次方程——而 \(y\) 又要靠上面那个三次方程求出来。

  真正把这几个公式一条条全写出来，能占满一整块黑板；所以数学家后来才问：
  五次能不能也有这样一条路？
  \liukong{1}
\end{example}

\noindent 上面这些公式，二次有、三次有、四次也有。五次方程，数学家找了三百多年。
十九世纪，伽罗瓦证明：五次方程\textbf{没有}求根公式。他用到的工具，就是群。

\begin{xiaojie}
  用自己的话，写下这门课里「群」到底是「什么」的名字。
  \liukong{4}
\end{xiaojie}

\noindent\textbf{全课终}：群不是新东西，它就是「绕圈」这件事的名字。而它，解开了三百年的谜。